% !PRINTED 565
\weberpage{565}
If the summation with respect to $t$ is carried out first,
then $st$ may be put in place of $t$, since $s$ is different from zero,
and hence $st$ and $t$ run at the same time through a complete system of residues
with respect to the modulus $n$. Thus
\weq{2}{(\varepsilon^\lambda, r)(\varepsilon^\mu, r) = \sum^{s} \sum^{t} r^{s(t+1)} \varepsilon^{(\lambda+\mu)\,\mathrm{ind}\,s} \varepsilon^{\mu\,\mathrm{ind}\,t}.}
We first dispose of the special case $\mu = -\lambda$, which
we obtain from (2) when we carry out the summation with respect to $s$
first:
\[ (\varepsilon^\lambda, r)(\varepsilon^{-\lambda}, r) = \sum^{t} \varepsilon^{-\lambda\,\mathrm{ind}\,t} \sum^{s} r^{s(t+1)}. \]
Herein
\[
\begin{aligned}
\sum^{s} r^{s(t+1)} &= -1, \quad \text{if } t = 1, 2 \ldots n-2\\
&= n-1, \quad \text{if } t = n-1,
\end{aligned}
\]
hence, since by \S.\,136 $\mathrm{ind}(n-1) = \frac{1}{2}(n-1)$ and
\[ \varepsilon^{\frac{n-1}{2}} = -1 \]
holds,
\[ (\varepsilon^\lambda, r)(\varepsilon^{-\lambda}, r) = -\sum_{1,\,n-1}^{t} \varepsilon^{-\lambda\,\mathrm{ind}\,t} + (-1)^\lambda n. \]
In the sum taken with respect to $t$, $\mathrm{ind}\,t$ runs through a complete
system of residues with respect to the modulus $n-1$, and hence, if $\lambda$ is
assumed \gesperrt{not divisible} by $n-1$,
\[ \sum_{1,\,n-1}^{t} \varepsilon^{-\lambda\,\mathrm{ind}\,t} = 0, \]
and therefore
\weq{3}{(\varepsilon^\lambda, r)(\varepsilon^{-\lambda}, r) = (-1)^\lambda n,}
while, if $\lambda$ is divisible by $n-1$, $(\varepsilon^\lambda, r) = (1, r) = -1$.

We now, then, treat the sum (2) further under the
assumption that $\lambda + \mu$ is not divisible by $(n-1)$.
Then $\sum^{s} \varepsilon^{(\lambda+\mu)\,\mathrm{ind}\,s} = 0$, and in the sum on the
right-hand side of (2) the terms corresponding to the value $t = n-1$
may be omitted. One thus obtains
\weq{4}{
\begin{aligned}
(\varepsilon^\lambda, r)(\varepsilon^\mu, r) &= \sum_{1,\,n-2}^{t} \varepsilon^{\mu\,\mathrm{ind}\,t} \sum_{1,\,n-1}^{s} \varepsilon^{(\lambda+\mu)\,\mathrm{ind}\,s} r^{s(t+1)}\\
&= \sum_{1,\,n-2}^{t} \varepsilon^{\mu\,\mathrm{ind}\,t - (\lambda+\mu)\,\mathrm{ind}(t+1)} \sum^{s} \varepsilon^{(\lambda+\mu)\,\mathrm{ind}\,s(t+1)} r^{s(t+1)}.
\end{aligned}
}
Now, for fixed $t$, $s(t+1)$ runs at the same time as $s$
through a complete system of residues with respect to the modulus $n$, and hence
%CONT>>
